% Part: sets-functions-relations
% Chapter: functions
% Section: functions-relations
% Hindi translation (hi-Deva-IN) of Open Logic commit 9620cc73f9c8e0ad003c514a5d3748f29611c4c0.

\documentclass[../../../include/open-logic-section]{subfiles}


\begin{document}

\olfileid[hi]{sfr}{fun}{rel}

\olsection{संबंधों के रूप में फलन}

\begin{explain}
जो फलन~$A$ के !!{element}ों को~$B$ के !!{element}ों पर प्रतिचित्रित करता है,
वह स्पष्टतः $A$ और~$B$ के बीच एक संबंध परिभाषित करता है; अर्थात वह संबंध जो
$x$ और $y$ के बीच तभी होता है जब $f(x) = y$। वास्तव में, यदि उदाहरणार्थ हम
गणित के आधारभूत घटकों को घटाने में रुचि रखते हों, तो फलन~$f$ को इस संबंध के
साथ, अर्थात युग्मों के एक समुच्चय के साथ, \emph{एक ही वस्तु मान} सकते हैं। इससे
फिर प्रश्न उठता है: कौन-से संबंध इस ढंग से फलन परिभाषित करते हैं?
\end{explain}

\begin{defn}[फलन का आलेख] मान लीजिए कि $f\colon A \to B$ एक फलन है।
फलन~$f$ का \emph{आलेख}, निम्न प्रकार परिभाषित संबंध
$R_f \subseteq A \times B$ है:
\[
R_f = \Setabs{\tuple{x,y}}{f(x) = y}.
\]
\end{defn}

\begin{explain}
विस्तारिता के कारण किसी फलन का आलेख अद्वितीय रूप से निर्धारित होता है। इसके
अतिरिक्त, (समुच्चयों पर) विस्तारिता फलनों के लिए विस्तारिता के निहित सिद्धांत
को तुरंत न्यायसंगत ठहराती है: यदि $f$ और~$g$ का प्रांत और सहप्रांत समान हो,
तो सभी मानों पर उनकी सहमति होने पर वे समान होते हैं।

इसी प्रकार, यदि कोई संबंध “फलनात्मक” है, तो वह किसी फलन का आलेख है।
\end{explain}

\begin{prop}\ollabel{prop:graph-function}
मान लीजिए कि $R \subseteq A \times B$ ऐसा है कि:
\begin{enumerate}
\item यदि $Rxy$ और $Rxz$, तो $y = z$; और
\item प्रत्येक $x \in A$ के लिए कोई $y \in B$ ऐसा है कि $\tuple{x,
y} \in R$।
\end{enumerate}
तब $R$ उस फलन $f\colon A \to B$ का आलेख है जिसे इस प्रकार परिभाषित किया गया
है कि $f(x) = y$ तभी जब $Rxy$।
\end{prop}

\begin{proof}
मान लीजिए कि कोई $y$ ऐसा है कि $Rxy$। यदि कोई अन्य $z \neq
y$ ऐसा होता कि
$Rxz$, तो~$R$ पर लगी शर्त का उल्लंघन होता। अतः यदि कोई $y$ ऐसा है कि $Rxy$,
तो यह $y$ अद्वितीय है; इसलिए $f$ सुपरिभाषित है। स्पष्टतः $R_f = R$।
\end{proof}

\begin{explain}
प्रत्येक फलन $f\colon A \to B$ का एक आलेख होता है, अर्थात $A
\times B$ पर वह संबंध जो $f(x) = y$ द्वारा परिभाषित है। दूसरी ओर,
\olref{prop:graph-function} में दिए गए गुणों वाला प्रत्येक संबंध~$R
\subseteq A \times B$, किसी फलन~$f \colon A \to
B$ का आलेख है। फलनों और
उनके आलेखों के बीच इस निकट संबंध के कारण हम किसी फलन को बस उसका आलेख मान सकते
हैं। दूसरे शब्दों में, फलनों को कुछ संबंधों के समान, अर्थात क्रमित समूहों के
कुछ समुच्चयों के समान, माना जा सकता है।
\oliflabeldef{sfr:rel:ref:sec}{फिर भी ध्यान दें कि इस “एक ही वस्तु मानने” का आशय
\olref[sfr][rel][ref]{sec} जैसा है: यह फलनों के तत्त्वमीमांसीय स्वरूप के बारे
में दावा नहीं, बल्कि यह अवलोकन है कि फलनों को कुछ समुच्चयों के रूप में
\emph{मानना} सुविधाजनक है। इस सुविधा का एक कारण यह है कि}{} अब हम फलनों पर
वैसी ही संक्रियाएँ करने पर विचार कर सकते हैं जैसी हमने संबंधों पर की थीं
(देखें \olref[sfr][rel][ops]{sec})। विशेषतः:
\end{explain}

\begin{defn}\ollabel{defn:funimage}
मान लीजिए कि $f \colon A \to B$ ऐसा फलन है जिसमें $C\subseteq A$।

फलन~$f$ का~$C$ पर \emph{प्रतिबंधन} वह फलन
$\funrestrictionto{f}{C}\colon C \to B$ है जिसे
$(\funrestrictionto{f}{C})(x) = f(x)$ द्वारा प्रत्येक $x \in C$ के लिए
परिभाषित किया जाता है। दूसरे
शब्दों में, $\funrestrictionto{f}{C} = \Setabs{\tuple{x, y} \in R_f}{x \in
C}$।

फलन~$f$ का~$C$ पर \emph{अनुप्रयोग} $\funimage{f}{C} =
\Setabs{f(x)}{x \in C}$ है। इसे~$C$ का~$f$ के अधीन \emph{प्रतिबिंब} भी कहते
हैं।
\end{defn}

\begin{explain}
इन परिभाषाओं से यह निष्कर्ष निकलता है कि $\ran{f} =
\funimage{f}{\dom{f}}$, किसी भी फलन~$f$ के लिए।
\oliflabeldef{sfr:rel:ops:sec}{फलनों को संबंध मानने और
\olref[sfr][rel][ops]{sec} की परिभाषाओं को देखते हुए, ये धारणाएँ ठीक वैसी ही
हैं जैसी अपेक्षित थीं। पर दो अन्य संक्रियाओं---प्रतिलोम और सापेक्ष
गुणनफलों---के लिए कुछ और विस्तार आवश्यक है। वह हम \olref[inv]{sec} और
\olref[cmp]{sec} में देंगे।}{}
\end{explain}

\end{document}
